% --- Archon physics formalization source begin ---
% archon:physics
% archon:covers PhyXMiniProblems/problem_phyx_mini_0256.lean
% archon:source-report reports/phyx_mini/problem_phyx_mini_0256.source.json

\chapter{Physics problem phyx\_mini\_0256}
\label{ch:PhyXMiniProblems_problem_phyx_mini_0256}

\paragraph{Problem source.}
\#\# Physical scenario

In figure, two springs are joined and connected to a block of mass \$0.245\textbackslash{} kg\$ that is set oscillating over a frictionless floor. The springs each have spring constant \$k = 6430\textbackslash{} N/m\$.

\#\# Question

What is the frequency of the oscillations?

\#\# Answer choices

- A: A: 16.8 Hz

- B: B: 17.6 Hz

- C: C: 18.2 Hz

- D: D: 19.4 Hz

\#\# Figure caption (auxiliary; use the image as primary evidence)

The image shows a mechanical system diagram. It features the following elements:

- A block labeled "m" on the left side, which represents a mass.
- Two springs connected in series, both labeled with "k," representing their spring constant.
- The springs are attached between the mass and a wall on the right.
- The entire setup is placed on a flat horizontal surface.
- The springs and mass are aligned horizontally, with the mass positioned to the left and the springs extending towards the right wall.

The image represents a classical physics problem involving mass and springs.

\#\# Dataset metadata

Wave Properties; Physical Model Grounding Reasoning, Multi-Formula Reasoning

\paragraph{Recorded answer/context.}
C: C: 18.2 Hz

\paragraph{Figure/image path.}
/root/proposal\_for\_physic/science-mango/phyx\_mini\_run/phyx\_data/test\_image/256.png

\paragraph{Formalization target.}
create a compiling Lean file with sorry bodies at `PhyXMiniProblems/problem\_phyx\_mini\_0256.lean`. The Lean declarations must preserve the physical quantities, dimensions or dimensional roles, figure labels, governing-law hypotheses, and final relation expressed by this problem.
Use Mathlib/Physlib names found through LeanExplore where available. If a physics API is missing, introduce faithful local abstractions rather than scalar placeholder aliases.

\begin{theorem}[Physics formalization target]
\label{thm:physics:phyx_mini_0256:target}
\lean{PhyXMiniProblems.ProblemPhyXMini0256.seriesSpringsFrequency_exact_and_matches_choiceC}
\uses{def:physics:phyx-mini-0256:phyxminiproblems-problemphyxmini0256-frequencyinhertz, def:physics:phyx-mini-0256:phyxminiproblems-problemphyxmini0256-seriesspringblocksetup, def:physics:phyx-mini-0256:phyxminiproblems-problemphyxmini0256-matchesproblemandfigure, def:physics:phyx-mini-0256:phyxminiproblems-problemphyxmini0256-hasphysicalparameters, def:physics:phyx-mini-0256:phyxminiproblems-problemphyxmini0256-satisfiesseriesspringoscillatormodel, def:physics:phyx-mini-0256:phyxminiproblems-problemphyxmini0256-matchesanswerchoice}
The assigned autoformalize agent should translate this physics problem into Lean declarations in the covered file, with theorem and lemma proof bodies written as `by sorry`.
\end{theorem}
\begin{proof}
This is an autoformalization task, not a proof task. Produce statements that can later be proved by the physics prover without weakening the physical model.
\end{proof}

% --- Archon declaration topology begin ---
\section{Declaration topology}
\begin{definition}[Mass Quantity]
  \label{def:physics:phyx-mini-0256:phyxminiproblems-problemphyxmini0256-massquantity}
  \lean{PhyXMiniProblems.ProblemPhyXMini0256.MassQuantity}
  A nonnegative physical mass
\end{definition}
\begin{proof}
  This is the declared model definition of Mass Quantity.
\end{proof}
\begin{definition}[Spring Stiffness Quantity]
  \label{def:physics:phyx-mini-0256:phyxminiproblems-problemphyxmini0256-springstiffnessquantity}
  \lean{PhyXMiniProblems.ProblemPhyXMini0256.SpringStiffnessQuantity}
  A nonnegative linear spring stiffness, whose SI unit is newtons per metre and whose physical dimension is mass per time squared
\end{definition}
\begin{proof}
  This is the declared model definition of Spring Stiffness Quantity.
\end{proof}
\begin{definition}[Frequency Quantity]
  \label{def:physics:phyx-mini-0256:phyxminiproblems-problemphyxmini0256-frequencyquantity}
  \lean{PhyXMiniProblems.ProblemPhyXMini0256.FrequencyQuantity}
  A nonnegative cyclic frequency, with inverse-time dimension
\end{definition}
\begin{proof}
  This is the declared model definition of Frequency Quantity.
\end{proof}
\begin{definition}[mass In Kilograms]
  \label{def:physics:phyx-mini-0256:phyxminiproblems-problemphyxmini0256-massinkilograms}
  \lean{PhyXMiniProblems.ProblemPhyXMini0256.massInKilograms}
  \uses{def:physics:phyx-mini-0256:phyxminiproblems-problemphyxmini0256-massquantity}
  Read a physical mass in SI kilograms
\end{definition}
\begin{proof}
  This is the declared model definition of mass In Kilograms.
\end{proof}
\begin{definition}[stiffness In Newtons Per Meter]
  \label{def:physics:phyx-mini-0256:phyxminiproblems-problemphyxmini0256-stiffnessinnewtonspermeter}
  \lean{PhyXMiniProblems.ProblemPhyXMini0256.stiffnessInNewtonsPerMeter}
  \uses{def:physics:phyx-mini-0256:phyxminiproblems-problemphyxmini0256-springstiffnessquantity}
  Read a spring stiffness in SI newtons per metre
\end{definition}
\begin{proof}
  This is the declared model definition of stiffness In Newtons Per Meter.
\end{proof}
\begin{definition}[frequency In Hertz]
  \label{def:physics:phyx-mini-0256:phyxminiproblems-problemphyxmini0256-frequencyinhertz}
  \lean{PhyXMiniProblems.ProblemPhyXMini0256.frequencyInHertz}
  \uses{def:physics:phyx-mini-0256:phyxminiproblems-problemphyxmini0256-frequencyquantity}
  Read a cyclic frequency in hertz, i.e. cycles per SI second
\end{definition}
\begin{proof}
  This is the declared model definition of frequency In Hertz.
\end{proof}
\begin{definition}[Spring Label]
  \label{def:physics:phyx-mini-0256:phyxminiproblems-problemphyxmini0256-springlabel}
  \lean{PhyXMiniProblems.ProblemPhyXMini0256.SpringLabel}
  The two springs, distinguished by their positions along the series chain
\end{definition}
\begin{proof}
  This is the declared model definition of Spring Label.
\end{proof}
\begin{definition}[Chain Endpoint]
  \label{def:physics:phyx-mini-0256:phyxminiproblems-problemphyxmini0256-chainendpoint}
  \lean{PhyXMiniProblems.ProblemPhyXMini0256.ChainEndpoint}
  Endpoints visible in the block-to-wall spring chain
\end{definition}
\begin{proof}
  This is the declared model definition of Chain Endpoint.
\end{proof}
\begin{definition}[Figure Label]
  \label{def:physics:phyx-mini-0256:phyxminiproblems-problemphyxmini0256-figurelabel}
  \lean{PhyXMiniProblems.ProblemPhyXMini0256.FigureLabel}
  Labels printed on the block and springs in the supplied figure
\end{definition}
\begin{proof}
  This is the declared model definition of Figure Label.
\end{proof}
\begin{definition}[Motion Axis]
  \label{def:physics:phyx-mini-0256:phyxminiproblems-problemphyxmini0256-motionaxis}
  \lean{PhyXMiniProblems.ProblemPhyXMini0256.MotionAxis}
  The horizontal direction of the block's allowed displacement
\end{definition}
\begin{proof}
  This is the declared model definition of Motion Axis.
\end{proof}
\begin{definition}[Floor Condition]
  \label{def:physics:phyx-mini-0256:phyxminiproblems-problemphyxmini0256-floorcondition}
  \lean{PhyXMiniProblems.ProblemPhyXMini0256.FloorCondition}
  The idealized contact condition specified in the problem
\end{definition}
\begin{proof}
  This is the declared model definition of Floor Condition.
\end{proof}
\begin{definition}[Two Series Spring Figure]
  \label{def:physics:phyx-mini-0256:phyxminiproblems-problemphyxmini0256-twoseriesspringfigure}
  \lean{PhyXMiniProblems.ProblemPhyXMini0256.TwoSeriesSpringFigure}
  \uses{def:physics:phyx-mini-0256:phyxminiproblems-problemphyxmini0256-springlabel, def:physics:phyx-mini-0256:phyxminiproblems-problemphyxmini0256-chainendpoint, def:physics:phyx-mini-0256:phyxminiproblems-problemphyxmini0256-figurelabel}
  Qualitative data transcribed from the primary image. The endpoint maps retain the order block--first spring--junction--second spring--wall, rather than collapsing the two drawn springs into a single glyph
\end{definition}
\begin{proof}
  This is the declared model definition of Two Series Spring Figure.
\end{proof}
\begin{definition}[Series Spring Block Setup]
  \label{def:physics:phyx-mini-0256:phyxminiproblems-problemphyxmini0256-seriesspringblocksetup}
  \lean{PhyXMiniProblems.ProblemPhyXMini0256.SeriesSpringBlockSetup}
  \uses{def:physics:phyx-mini-0256:phyxminiproblems-problemphyxmini0256-massquantity, def:physics:phyx-mini-0256:phyxminiproblems-problemphyxmini0256-springstiffnessquantity, def:physics:phyx-mini-0256:phyxminiproblems-problemphyxmini0256-frequencyquantity, def:physics:phyx-mini-0256:phyxminiproblems-problemphyxmini0256-springlabel, def:physics:phyx-mini-0256:phyxminiproblems-problemphyxmini0256-motionaxis, def:physics:phyx-mini-0256:phyxminiproblems-problemphyxmini0256-floorcondition, def:physics:phyx-mini-0256:phyxminiproblems-problemphyxmini0256-twoseriesspringfigure}
  The independent physical quantities and observables of the apparatus. oscillationFrequency and effectiveSeriesStiffness are unknown physical quantities. The scalar Physlib oscillator is linked to them only by the general laws below; none of these fields defines the requested answer
\end{definition}
\begin{proof}
  This is the declared model definition of Series Spring Block Setup.
\end{proof}
\begin{definition}[Matches Problem And Figure]
  \label{def:physics:phyx-mini-0256:phyxminiproblems-problemphyxmini0256-matchesproblemandfigure}
  \lean{PhyXMiniProblems.ProblemPhyXMini0256.MatchesProblemAndFigure}
  \uses{def:physics:phyx-mini-0256:phyxminiproblems-problemphyxmini0256-massinkilograms, def:physics:phyx-mini-0256:phyxminiproblems-problemphyxmini0256-stiffnessinnewtonspermeter, def:physics:phyx-mini-0256:phyxminiproblems-problemphyxmini0256-springlabel, def:physics:phyx-mini-0256:phyxminiproblems-problemphyxmini0256-seriesspringblocksetup}
  Problem-text calibrations and qualitative evidence from the supplied image. This fixes the two component stiffnesses, but neither the effective stiffness nor the requested oscillation frequency
\end{definition}
\begin{proof}
  This is the declared model definition of Matches Problem And Figure.
\end{proof}
\begin{definition}[Has Physical Parameters]
  \label{def:physics:phyx-mini-0256:phyxminiproblems-problemphyxmini0256-hasphysicalparameters}
  \lean{PhyXMiniProblems.ProblemPhyXMini0256.HasPhysicalParameters}
  \uses{def:physics:phyx-mini-0256:phyxminiproblems-problemphyxmini0256-massinkilograms, def:physics:phyx-mini-0256:phyxminiproblems-problemphyxmini0256-stiffnessinnewtonspermeter, def:physics:phyx-mini-0256:phyxminiproblems-problemphyxmini0256-springlabel, def:physics:phyx-mini-0256:phyxminiproblems-problemphyxmini0256-seriesspringblocksetup}
  Positivity of the mass, component stiffnesses, and effective stiffness
\end{definition}
\begin{proof}
  This is the declared model definition of Has Physical Parameters.
\end{proof}
\begin{definition}[Satisfies Series Spring Oscillator Model]
  \label{def:physics:phyx-mini-0256:phyxminiproblems-problemphyxmini0256-satisfiesseriesspringoscillatormodel}
  \lean{PhyXMiniProblems.ProblemPhyXMini0256.SatisfiesSeriesSpringOscillatorModel}
  \uses{def:physics:phyx-mini-0256:phyxminiproblems-problemphyxmini0256-massinkilograms, def:physics:phyx-mini-0256:phyxminiproblems-problemphyxmini0256-stiffnessinnewtonspermeter, def:physics:phyx-mini-0256:phyxminiproblems-problemphyxmini0256-frequencyinhertz, def:physics:phyx-mini-0256:phyxminiproblems-problemphyxmini0256-seriesspringblocksetup}
  Governing small-oscillation laws For two linear springs in series, equal force and additive extension give k\_eff * (k₁ + k₂) = k₁ * k₂. The equivalent Physlib oscillator has the same block mass and effective stiffness. Physlib supplies HarmonicOscillator.ω = sqrt (k / m); cyclic frequency is related to angular frequency by ω = 2 π f. These are general modeling laws and contain no numerical frequency or answer choice for this problem
\end{definition}
\begin{proof}
  This is the declared model definition of Satisfies Series Spring Oscillator Model.
\end{proof}
\begin{definition}[Answer Choice]
  \label{def:physics:phyx-mini-0256:phyxminiproblems-problemphyxmini0256-answerchoice}
  \lean{PhyXMiniProblems.ProblemPhyXMini0256.AnswerChoice}
  Labels of the four answers printed with the problem
\end{definition}
\begin{proof}
  This is the declared model definition of Answer Choice.
\end{proof}
\begin{definition}[hertz]
  \label{def:physics:phyx-mini-0256:phyxminiproblems-problemphyxmini0256-answerchoice-hertz}
  \lean{PhyXMiniProblems.ProblemPhyXMini0256.AnswerChoice.hertz}
  \uses{def:physics:phyx-mini-0256:phyxminiproblems-problemphyxmini0256-answerchoice}
  Frequency in hertz printed beside each answer label
\end{definition}
\begin{proof}
  This is the declared model definition of hertz.
\end{proof}
\begin{definition}[Matches Answer Choice]
  \label{def:physics:phyx-mini-0256:phyxminiproblems-problemphyxmini0256-matchesanswerchoice}
  \lean{PhyXMiniProblems.ProblemPhyXMini0256.MatchesAnswerChoice}
  \uses{def:physics:phyx-mini-0256:phyxminiproblems-problemphyxmini0256-frequencyquantity, def:physics:phyx-mini-0256:phyxminiproblems-problemphyxmini0256-frequencyinhertz, def:physics:phyx-mini-0256:phyxminiproblems-problemphyxmini0256-answerchoice}
  Agreement with a displayed one-decimal-place answer, using the usual nearest tenth interval of one twentieth hertz on either side
\end{definition}
\begin{proof}
  This is the declared model definition of Matches Answer Choice.
\end{proof}
% --- Archon declaration topology end ---
% --- Archon physics formalization source end ---
